\chapter*{Avant-propos et mode d'emploi}
\addcontentsline{toc}{chapter}{Avant-propos et mode d'emploi}
\markboth{Avant-propos}{Avant-propos}

Ce rapport poursuit trois objectifs : \textbf{documenter} intégralement le modèle économique et social de la France
construit dans le projet « Présidentielle » ; l'\textbf{utiliser} pour comprendre, par des expériences chiffrées, comment
les politiques publiques se propagent dans l'économie et la société ; en déduire des \textbf{stratégies de redressement}
et en mesurer les arbitrages. Il s'adresse à deux lecteurs à la fois : le lecteur sans formation économique, pour qui
des encadrés pédagogiques et un glossaire (annexe~\ref{ann:glossaire}) expliquent les notions ; l'économiste, qui
trouvera les équations telles qu'elles sont codées, les coefficients estimés avec leurs écarts-types, les sources et
les tests, de quoi critiquer, vérifier et reproduire chaque résultat.

\paragraph{Ce que ce rapport n'est pas.} Ce n'est ni une prévision, ni un programme politique. Toutes les simulations
comparent deux avenirs hypothétiques (avec et sans une mesure) et dépendent d'hypothèses exposées à chaque étape. Le
modèle est un outil de raisonnement chiffré : il rend explicites des chaînes de causalité et en donne des ordres de
grandeur. Il ne démontre pas qu'une mesure « marche » : il montre ce qu'elle produit \emph{si} les comportements du passé
restent valables et \emph{si} les hypothèses retenues sont justes.

\paragraph{Règles suivies.}
\begin{itemize}
\item Toutes les données sont ouvertes ; les sources sont citées au niveau des affirmations.
\item Chaque paramètre a un niveau de preuve (relation comptable, donnée observée, estimé, calibré, transposé de la
littérature, hypothèse) ; chaque résultat un degré de robustesse (\fiab{robuste}, \fiab{plutôt robuste}, \fiab{fragile},
\fiab{non établi}).
\item Une corrélation passée n'est jamais présentée comme un effet causal établi : les coefficients estimés sont des
régularités historiques, dont la validité pour des réformes de grande ampleur n'est pas garantie.
\item Les résultats du modèle (parties I à III) sont séparés des recommandations (partie IV), elles-mêmes conditionnelles
aux choix de pondération entre objectifs, qui sont des choix politiques.
\item Toutes les mesures passent par les mêmes équations, quelle que soit leur origine politique.
\item Tous les chiffres proviennent de simulations réellement exécutées avec le code du dépôt ; les fichiers de résultats
sont fournis.
\end{itemize}

\paragraph{Plan.} La \textbf{partie I} décrit l'architecture du modèle (chapitre 1), chacun de ses sous-modèles avec ses
équations, sa calibration et ses limites (chapitre 2), puis en fait l'audit (chapitre 3). La \textbf{partie II} présente
la situation de référence comparée aux statistiques (chapitre 4), huit expériences qui isolent les mécanismes
(chapitre 5) et des expériences complémentaires sur les interactions (chapitre 6). La \textbf{partie III} recherche des
programmes de redressement par optimisation et les soumet à des tests de robustesse (chapitre 7). La \textbf{partie IV}
répond aux huit questions de la mission et propose un programme hiérarchisé (chapitre 8).

\section*{Résumé des principaux enseignements}

\begin{enumerate}
\item \textbf{La trajectoire « sans mesure » n'est pas soutenable.} Avec un taux d'intérêt de marché de 4,9~\% et une
croissance faible (1,1 à 1,2~\% par an jusqu'en 2032, moins ensuite), le déficit public passerait de 5,1~\% du PIB en 2025 à
9~\% en 2032 et la dette de 116 à 142~\% du PIB. La charge d'intérêts en est le premier moteur (de 2,5 à 5,1 points de PIB),
devant le vieillissement. \fiab{robuste} pour le sens, \fiab{fragile} pour les niveaux après 2035.
\item \textbf{Toute consolidation coûte de l'activité à court terme, mais pas autant selon l'instrument.} Par point de
déficit réduit, une économie sur la production publique (fonctionnement, emplois) coûte environ 2~\% de PIB la première
année ; une hausse de TVA, une économie sur les transferts ou une hausse d'impôt sur le patrimoine, 0,3 à 0,7~\%. Les
effets sur l'activité s'estompent en cinq à dix ans ; les gains sur la dette s'accumulent. \fiab{plutôt robuste}.
\item \textbf{Une hausse d'impôt rapporte d'un quart à un tiers de moins que son rendement statique} la première année,
à cause de l'activité perdue ; aucune hausse d'impôt testée ne réduit les recettes, sauf la tranche supérieure de l'IR si
les contribuables aisés réagissent fortement. \fiab{robuste} pour le signe.
\item \textbf{Les mesures qui réduisent le plus la pauvreté par euro sont des prestations ciblées} (prime d'activité, RSA,
minimum vieillesse). Elles soutiennent l'activité à court terme. Les impôts sur la consommation (TVA, taxe carbone) sont
ceux qui pèsent le plus sur les ménages modestes. \fiab{plutôt robuste} (effets statiques).
\item \textbf{Une hausse générale des salaires ne crée ni ne détruit beaucoup d'activité} : elle redistribue vers les
salariés en emploi, au détriment des revenus fixes et, à moyen terme, de l'emploi. Elle n'augmente pas la pauvreté par
effet de seuil (le revenu disponible croît toujours avec le salaire), mais les taux marginaux effectifs atteignent 55 à
82~\% entre 1 et 2 SMIC. \fiab{plutôt robuste}.
\item \textbf{Seul l'investissement public a, parmi les dépenses, un effet durable sur la production} dans le modèle ;
le report de l'âge de départ est le levier de production le plus puissant à long terme, mais il augmente d'abord le
chômage et ses effets dépendent d'une hypothèse de bouclage non estimée. \fiab{fragile}.
\item \textbf{Les mesures interagissent} : relocalisation industrielle et droits de douane se neutralisent (la riposte et
la hausse des coûts annulent le bénéfice), alors que l'investissement public et le report de l'âge se complètent.
\item \textbf{Aucun programme testé ne stabilise la dette si les taux restent proches de 5~\%.} L'effort nécessaire
(environ 6 points de PIB) dépasse ce que les leviers disponibles peuvent apporter sans récession majeure : un effort
maximal (scénario D) coûte deux années de récession et 6~\% de pouvoir d'achat durablement. La soutenabilité dépend
d'abord des taux d'intérêt et de la croissance.
\item \textbf{Le programme recommandé} (chapitre~\ref{ch:synthese}) combine une consolidation par les instruments les
moins coûteux en activité, une protection ciblée des ménages modestes, l'investissement public et un report progressif
de l'âge de départ à 65 ans. Il réduit le déficit, la pauvreté et la dette par rapport à la référence et augmente la
production à long terme, sans la stabiliser à lui seul : c'est un socle, pas une solution complète.
\item \textbf{Les principales faiblesses du modèle} sont l'absence d'estimation du bouclage de moyen terme, une
représentation de l'offre (productivité, démographie) qui donne des résultats paradoxaux pour les chocs d'offre,
l'absence de bilans financiers et d'anticipations, des dépenses publiques indistinctes, et une microsimulation statique
qui représente mal le haut de la distribution (chapitre~\ref{ch:audit}).
\end{enumerate}
